\section{从坐标环得到切空间}\label{sec:intrinsic-tangent-spaces}

本节的目标是完全用坐标环 $K[X]$ 重新解释切空间 $T_{a}X$. 这一内在描述的一个有用推论是, 切空间在同构下保持不变. 换句话说, 若 $F:X\to Y$ 为同构, 且 $a\in X$, 则 $T_{a}X$ 和 $T_{F(a)}Y$ 是同构的向量空间.

为了开始给出切空间的内在描述, 首先需要一种内在的方式来理解点 $a\in X$. 给定 $a\in X$, 定义
\[
  I_{a}=\{F\in K[X]\mid F(a)=0\}\subseteq K[X].
\]
由定义可见, $I_{a}$ 是 $K[X]$ 的子集. 不止如此, $I_{a}$ 是 $K[X]$ 的极大理想, 建议读者在 Exercise~\ref{ex:point-ideal-is-maximal-in-coordinate-ring} 中验证. 注意 $I_{a}$ 是内在定义的, 不涉及 $X$ 所处的环境仿射空间.

为更具体地描述 $I_{a}$, 考虑 $K[X]$ 的一种 (外在的) 商环描述. 更具体地说, 假设 $X\subseteq\mathbb{A}^{n}$, 于是可写 $K[X]=K[x_{1},\ldots,x_{n}]/\mathcal{I}(X)$ 和 $a=(a_{1},\ldots,a_{n})\in X$. 对任意 $F\in K[X]$, 可对某个 $f\in K[x_{1},\ldots,x_{n}]$ 写作 $F=[f]$. 那么 $F\in I_{a}$ 当且仅当 $f(a)=0$. 因为每个在 $a$ 处取值为零的多项式都能写为
\[
  f=(x_{1}-a_{1})f_{1}+\cdots+(x_{n}-a_{n})f_{n},
\]
其中 $f_{1},\ldots,f_{n}\in K[x_{1},\ldots,x_{n}]$ 为某些多项式, 所以
\begin{equation}\label{eq:point-ideal-from-coordinate-generators}
  I_{a}=\langle[x_{1}-a_{1}],\ldots,[x_{n}-a_{n}]\rangle
    \subseteq\frac{K[x_{1},\ldots,x_{n}]}{\mathcal{I}(X)}.
\end{equation}
注意, $I_{a}$ 可以自然地看成向量空间: 两个在 $a$ 处取值为零的函数之和在 $a$ 处取值为零, 一个在 $a$ 处取值为零的函数乘以标量后也在 $a$ 处取值为零.

现在考虑理想 $I_{a}^{2}=I_{a}\cdot I_{a}$, 即 $I_{a}$ 与自身的理想乘积. 利用等式 \eqref{eq:point-ideal-from-coordinate-generators} 中 $I_{a}$ 的生成元, 可得
\begin{equation}\label{eq:square-of-point-ideal}
  I_{a}^{2}=\langle[(x_{i}-a_{i})(x_{j}-a_{j})]\mid1\leq i,j\leq n\rangle\subseteq I_{a}.
\end{equation}
注意 $I_{a}^{2}$ 是 $I_{a}$ 的向量子空间, 所以可以考虑商向量空间 $I_{a}/I_{a}^{2}$. 用等式 \eqref{eq:point-ideal-from-coordinate-generators} 中的生成元来看, 该商是有限维的, 由线性函数 $[x_{i}-a_{i}]$ 张成 (Exercise~\ref{ex:point-ideal-quotient-spanned-by-coordinate-linear-functions}):
\[
  I_{a}/I_{a}^{2}=K\{[x_{1}-a_{1}],\ldots,[x_{n}-a_{n}]\}.
\]
本节的主要结论如下, 它将有限维向量空间 $I_{a}/I_{a}^{2}$ 解释为切空间 $T_{a}X$ 的对偶.

\begin{Theorem}{$I_{a}/I_{a}^{2}=(T_{a}X)^{\vee}$}\leavevmode\label{thm:point-ideal-square-quotient-is-dual-of-tangent-space}
设 $X\subseteq\mathbb{A}^{n}$ 为仿射簇, 且 $a\in X$. 则有典范的向量空间同构
\[
  I_{a}/I_{a}^{2}=(T_{a}X)^{\vee}.
\]
\end{Theorem}

在给出证明之前, 先考察几个熟悉的切空间例子, 以便理解 Theorem~\ref{thm:point-ideal-square-quotient-is-dual-of-tangent-space} 中的新概念.

\begin{Example}{抛物线的切空间}\leavevmode\label{ex:tangent-space-of-parabola}
考虑抛物线 $X=\mathcal{V}(y-x^{2})\subseteq\mathbb{A}^{2}$, 令 $a=(0,0)$. 在 Example~\ref{ex:linearizations-of-parabola} 中, 我们算出 $X$ 有一维线性化 $L_{a}X$, 因而切空间 $T_{a}X$ 是一维的. 由 Theorem~\ref{thm:point-ideal-square-quotient-is-dual-of-tangent-space}, 可知 $I_{a}/I_{a}^{2}$ 应为一维向量空间. 为直接验证这一点, 注意
\[
  \begin{aligned}
    I_{a}&=\langle[x],[y]\rangle\\
      &=K\{[x],[y],[x^{2}],[xy],[y^{2}],[x^{3}],[x^{2}y],[xy^{2}],[y^{3}],\ldots\}\subseteq K[X]
  \end{aligned}
\]
且
\[
  \begin{aligned}
    I_{a}^{2}&=\langle[x^{2}],[xy],[y^{2}]\rangle\\
      &=K\{[x^{2}],[xy],[y^{2}],[x^{3}],[x^{2}y],[xy^{2}],[y^{3}],\ldots\}\subseteq K[X].
  \end{aligned}
\]
注意到在 $K[X]$ 中 $[y]=[x^{2}]$, 可知除 $[x]$ 之外, $I_{a}$ 的上述每个生成元也出现在 $I_{a}^{2}$ 中. 所以 $I_{a}/I_{a}^{2}=K\{[x]\}$, 这如预期是 $K$ 上的一维向量空间. 此外, 函数 $[x]\in K[X]$ 在线性化 $L_{a}X=\mathcal{V}(y)\subseteq\mathbb{A}^{2}$ 上给出线性函数. 将 $L_{a}X=\{(a,0)\mid a\in K\}$ 与 $T_{a}X=K$ 认同后, 就可以自然地将向量空间 $I_{a}/I_{a}^{2}=K\{[x]\}$ 与对偶空间 $(T_{a}X)^{\vee}$ 认同.
\end{Example}

\begin{Example}{尖点曲线的切空间}\leavevmode\label{ex:tangent-space-of-cusp}
如 Example~\ref{ex:linearizations-of-cusp}, 考虑尖点曲线 $X=\mathcal{V}(x^{2}-y^{3})\subseteq\mathbb{A}^{2}$, 令 $a=(0,0)$. 在 Exercise~\ref{ex:linearizations-of-cusp-at-origin-and-elsewhere} 中, 建议读者证明 $X$ 有二维线性化 $L_{a}X=\mathbb{A}^{2}$, 从而有二维切空间 $T_{a}X=K^{2}$. 由 Theorem~\ref{thm:point-ideal-square-quotient-is-dual-of-tangent-space} 可知 $I_{a}/I_{a}^{2}$ 应是二维向量空间.

为直接验证这一预期, 完全按上一例计算 $I_{a}$ 和 $I_{a}^{2}$. 但与上一例不同的是, $[x]$ 和 $[y]$ 在商 $I_{a}/I_{a}^{2}$ 中仍然线性独立. 所以 $I_{a}/I_{a}^{2}=K\{[x],[y]\}$ 是二维向量空间. 将 $L_{a}X=\mathbb{A}^{2}$ 与 $T_{a}X=K^{2}$ 认同后, 可以自然地将向量空间 $I_{a}/I_{a}^{2}=K\{[x],[y]\}$ 看作对偶空间 $(T_{a}X)^{\vee}$.
\end{Example}

\begin{Proof}\textbf{of Theorem~\ref{thm:point-ideal-square-quotient-is-dual-of-tangent-space}.}\quad
我们定义典范的满射线性映射
\[
  \phi:I_{a}\to(T_{a}X)^{\vee}
\]
并证明 $\ker(\phi)=I_{a}^{2}$. 为定义 $\phi$, 设 $F\in I_{a}$, 并对某个 $f\in K[x_{1},\ldots,x_{n}]$ 写作 $F=[f]$. 定义 $\phi(F)$ 为线性映射
\[
  \begin{aligned}
    \phi(F):T_{a}X&\to K,\\
    v&\mapsto\nabla f(a)\cdot v.
  \end{aligned}
\]
当然, 我们应当顾虑这一定义是否依赖于代表元 $f$ 的选择. 然而, 由 Proposition~\ref{prop:gradient-characterization-of-tangent-spaces} 可知 (Exercise~\ref{ex:tangent-functional-independent-of-representative})
\[
  [f]=[g]\quad\Longrightarrow\quad
  \nabla f(a)\cdot v=\nabla g(a)\cdot v\qquad\text{对所有 }v\in T_{a}X,
\]
所以 $\phi(F)\in(T_{a}X)^{\vee}$ 与代表元 $f$ 的选择无关. 验证 $\phi$ 的线性性, 要求对任意 $F,G\in K[X]$ 和 $r\in K$ 都有 $\phi(F+rG)=\phi(F)+r\phi(G)$. 这由求导的线性性得到 (Exercise~\ref{ex:map-from-point-ideal-to-tangent-dual-is-linear}). 所以只剩证明 $\phi$ 是满射且 $\ker(\phi)=I_{a}^{2}$.

为证明 $\phi$ 是满射, 设 $\rho:T_{a}X\to K$ 为线性映射, 并将它扩张为线性映射 $\overline{\rho}:K^{n}\to K$. 定义线性多项式
\[
  f=\overline{\rho}(e_{1})(x_{1}-a_{1})+\cdots+\overline{\rho}(e_{n})(x_{n}-a_{n})
    \in K[x_{1},\ldots,x_{n}],
\]
其中 $e_{1},\ldots,e_{n}$ 是 $K^{n}$ 的标准基向量, 并注意 $[f]\in I_{a}$. 我们断言 $\phi([f])=\rho$. 将定义展开, 可见对所有 $i=1,\ldots,n$, 有 $\nabla f(a)\cdot e_{i}=\overline{\rho}(e_{i})$. 由于 $T_{a}X\subseteq K^{n}$ 上的线性映射由任意一个扩张在标准基向量上的取值唯一确定, 对所有 $v\in T_{a}X$ 都有 $\nabla f(a)\cdot v=\rho(v)$, 所以 $\phi([f])=\rho$. 因此 $\phi$ 是满射.

为证明 $\ker(\phi)=I_{a}^{2}$, 先设 $F\in I_{a}^{2}$. 由等式 \eqref{eq:square-of-point-ideal}, 有 $F=[f]$, 其中 $f$ 的形式为
\[
  f=\sum_{i,j=1}^{n}(x_{i}-a_{i})(x_{j}-a_{j})f_{i,j}
  \qquad\text{其中 }f_{i,j}\in K[x_{1},\ldots,x_{n}].
\]
利用乘积法则求 $f$ 的导数, 得 $\nabla f(a)=0$, 从而 $F\in\ker(\phi)$. 因此 $I_{a}^{2}\subseteq\ker(\phi)$.

为证明反向包含, 设 $F\in\ker(\phi)$, 并对某个 $f\in K[x_{1},\ldots,x_{n}]$ 写作 $F=[f]$. 假设 $F\in\ker(\phi)$ 意味着
\[
  \nabla f(a)\cdot v=0\qquad\text{对所有 }v\in T_{a}X.
\]
选取生成元 $\mathcal{I}(X)=\langle g_{1},\ldots,g_{m}\rangle$, 则由 Proposition~\ref{prop:gradient-characterization-of-tangent-spaces} 和线性代数的论证 (Exercise~\ref{ex:annihilator-of-kernel-is-row-span}) 得
\[
  \nabla f(a)=\sum_{i=1}^{m}a_{i}\nabla g_{i}(a),
\]
其中 $a_{1},\ldots,a_{m}\in K$ 为某些标量. 由于梯度对多项式线性作用, 且 $\mathcal{I}(X)$ 对线性组合封闭, 可见
\[
  \nabla f(a)=\nabla g(a)\qquad\text{其中 }g=\sum_{i=1}^{m}a_{i}g_{i}\in\mathcal{I}(X).
\]
利用 $F=[f]$ 属于 $\phi$ 的定义域从而 $[f]\in I_{a}$, 可以写成
\[
  f=\sum_{i=1}^{n}b_{i}(x_{i}-a_{i})
    +\sum_{i,j=1}^{n}f_{i,j}(x_{i}-a_{i})(x_{j}-a_{j}).
\]
对 $g$ 写出类似的表达式, 则等式 $\nabla f(a)=(b_{1},\ldots,b_{n})=\nabla g(a)$ 蕴含
\[
  f-g=\sum_{i,j=1}^{n}(f_{i,j}-g_{i,j})(x_{i}-a_{i})(x_{j}-a_{j}).
\]
由于 $g\in\mathcal{I}(X)$, 有 $[f]=[f-g]$, 从上面的表达式得到
\[
  F=[f-g]\in\langle[(x_{i}-a_{i})(x_{j}-a_{j})]\mid1\leq i,j\leq n\rangle=I_{a}^{2},
\]
证明完毕.
\end{Proof}

对 Theorem~\ref{thm:point-ideal-square-quotient-is-dual-of-tangent-space} 中的两边取对偶, 并应用 Theorem~\ref{thm:vector-space-canonically-isomorphic-to-double-dual}, 得到切空间的以下内在刻画.

\begin{Corollary}{$T_{a}X=(I_{a}/I_{a}^{2})^{\vee}$}\leavevmode\label{cor:tangent-space-as-dual-of-point-ideal-quotient}
设 $X\subseteq\mathbb{A}^{n}$ 为仿射簇, 且 $a\in X$. 则有典范的向量空间同构
\[
  T_{a}X=(I_{a}/I_{a}^{2})^{\vee}.
\]
\end{Corollary}

Theorem~\ref{thm:point-ideal-square-quotient-is-dual-of-tangent-space} 和 Corollary~\ref{cor:tangent-space-as-dual-of-point-ideal-quotient} 的主要收获是, 现在研究仿射簇的切空间无需提及该簇所处的仿射空间. 更具体地说, 我们已经看到, 任意一点 $a\in X$ 对应于坐标环 $K[X]$ 中的一个极大理想 $I_{a}$, 现在只须将切空间 $T_{a}X$ 抽象地看成向量空间 $(I_{a}/I_{a}^{2})^{\vee}$ 来研究. 正如你可能猜到的那样, 切空间确实是内在的, 即它们在同构下保持不变. 本节最后明确说明切空间的这种内在性.

\begin{Corollary}{切空间是内在的}\leavevmode\label{cor:tangent-spaces-invariant-under-isomorphism}
设 $X$ 和 $Y$ 为仿射簇, $F:X\to Y$ 为同构, 且 $a\in X$. 则 $F$ 诱导向量空间同构
\[
  T_{a}X\cong T_{F(a)}Y.
\]
\end{Corollary}
\begin{Proof}
建议读者在 Exercise~\ref{ex:pullback-identifies-maximal-point-ideals} 中验证: 拉回同构 $F^{*}:K[Y]\to K[X]$ 将极大理想 $I_{F(a)}$ 与 $I_{a}$ 认同:
\[
  F^{*}(I_{F(a)})=I_{a}.
\]
由此得
\[
  F^{*}(I_{F(a)}^{2})=I_{a}^{2}.
\]
因而 $F^{*}$ 诱导向量空间同构 $I_{F(a)}/I_{F(a)}^{2}\cong I_{a}/I_{a}^{2}$, 取对偶即得同构 $T_{a}X\cong T_{F(a)}Y$.
\end{Proof}

\BookExerciseHeading
\begin{BookExercises}
  \item\label{ex:point-ideal-is-maximal-in-coordinate-ring} 设 $X$ 为仿射簇, $a\in X$ 为一点. 证明
  \[
    I_{a}=\{F\in K[X]\mid F(a)=0\}
  \]
  是 $K[X]$ 的极大理想.
  \item\label{ex:point-ideal-quotient-spanned-by-coordinate-linear-functions} 设 $X\subseteq\mathbb{A}^{n}$ 为仿射簇, 且 $a\in X$. 证明
  \[
    I_{a}/I_{a}^{2}=K\{[x_{1}-a_{1}],\ldots,[x_{n}-a_{n}]\}.
  \]
  \item\label{ex:tangent-functional-independent-of-representative} 设 $X\subseteq\mathbb{A}^{n}$ 为仿射簇, 且 $a\in X$. 证明若 $[f]=[g]\in K[X]$, 则对所有 $v\in T_{a}X$, 有 $\nabla f(a)\cdot v=\nabla g(a)\cdot v$.
  \item\label{ex:map-from-point-ideal-to-tangent-dual-is-linear} 设 $\phi$ 如 Theorem~\ref{thm:point-ideal-square-quotient-is-dual-of-tangent-space} 的证明中所定义. 对任意 $F,G\in K[X]$ 和 $r\in K$, 利用求导的线性性证明
  \[
    \phi(F+rG)=\phi(F)+r\phi(G).
  \]
  \item\label{ex:annihilator-of-kernel-is-row-span} 设 $v_{1},\ldots,v_{m}\in K^{n}$, 且 $V\subseteq K^{n}$ 是由下式给出的线性子空间:
  \[
    V=\{v\in K^{n}\mid v_{i}\cdot v=0\text{ 对所有 }i=1,\ldots,m\}.
  \]
  假设 $w\in K^{n}$ 对所有 $v\in V$ 都满足 $w\cdot v=0$. 证明 $w$ 属于 $v_{1},\ldots,v_{m}$ 的张成空间. (提示: 考虑以 $v_{1},\ldots,v_{m}$ 为各行的矩阵 $M$, 以及把行 $w$ 添加到 $M$ 得到的矩阵 $M'$. 用秩--零度论证证明 $\operatorname{rk}(M)=\operatorname{rk}(M')$.)
  \item\label{ex:pullback-identifies-maximal-point-ideals} 设 $X\subseteq\mathbb{A}^{m}$ 和 $Y\subseteq\mathbb{A}^{n}$ 为仿射簇, $F:X\to Y$ 为同构, 且 $a\in X$. 证明 $F^{*}:K[Y]\to K[X]$ 将极大理想 $I_{F(a)}$ 与 $I_{a}$ 认同:
  \[
    F^{*}(I_{F(a)})=I_{a}.
  \]
\end{BookExercises}
